% GENERATED FILE. Edit canonical lessons, then run npm run generate:books.
% canonical-source-hash: fnv1a64:d3aac0e2051b819f
% canonical-lessons: ML-C188-maykkuka, ML-C188-onakkuka, ML-C188-ophakkuka, ML-C188-prasiddhikarikkuka, ML-C188-nannakkuka

\chapter{To erase, To switch on, To switch off, To publish, To repair}
\label{ch:ml-to-erase-to-switch-on-to-switch-off-to-publish-to-repair}
\hlchaptermodality{188}

\begin{chapteropening}
\textbf{By the end of this chapter:} \emph{I can say to erase, to switch on, to switch off, to publish and to repair.}

Complete the last lesson of chapter 188: I can say to erase, to switch on, to switch off, to publish and to repair.
\end{chapteropening}

\section[\texorpdfstring{\ml{മായ്ക്കുക} (māykkuka)}{മായ്ക്കുക (māykkuka)}]{\ml{മായ്ക്കുക} (māykkuka) — to erase, to delete}

\label{lesson:ML-C188-maykkuka}



\begin{quote}
Before the new one: say the Malayalam for to bark, then the Malayalam for to rise.
\end{quote}

\subsection*{You'll want to know: \ml{മായ്ക്കുക}}
\textbf{\ml{മായ്ക്കുക}} — \emph{māykkuka} — \textquotedblleft{}to erase, to delete\textquotedblright{}.

\begin{grammarlens}[title={What you've built}]
One more word for this chapter.
\end{grammarlens}

\subsection*{Your turn}
\begin{itemize}
\raggedright
  \item \emph{Say it:} \emph{māykkuka}
  \item \emph{Say it:} \emph{māykkuka}, once more
  \item \emph{Say it:} say \emph{māykkuka}
  \item \emph{Recall:} say the Malayalam for to bark, then the Malayalam for to rise, then say \emph{māykkuka} again
\end{itemize}

\begin{tcolorbox}[breakable,colback=teal!4,colframe=teal!35!black,title={Before you move on}]
What does \ml{മായ്ക്കുക} mean? (\textquotedblleft{}To erase, to delete\textquotedblright{}.) Say it once more.
\end{tcolorbox}
\section[\texorpdfstring{\ml{ഓണാക്കുക} (ōṇākkuka)}{ഓണാക്കുക (ōṇākkuka)}]{\ml{ഓണാക്കുക} (ōṇākkuka) — to switch on}

\label{lesson:ML-C188-onakkuka}



\begin{quote}
Before the new one: say the Malayalam for to rise, then the Malayalam for to erase.
\end{quote}

\subsection*{You'll want to know: \ml{ഓണാക്കുക}}
\textbf{\ml{ഓണാക്കുക}} — \emph{ōṇākkuka} — \textquotedblleft{}to switch on\textquotedblright{}.

\begin{grammarlens}[title={What you've built}]
One more word for this chapter.
\end{grammarlens}

\subsection*{Your turn}
\begin{itemize}
\raggedright
  \item \emph{Say it:} \emph{ōṇākkuka}
  \item \emph{Say it:} \emph{ōṇākkuka}, once more
  \item \emph{Say it:} say \emph{ōṇākkuka}
  \item \emph{Recall:} say the Malayalam for to rise, then the Malayalam for to erase, then say \emph{ōṇākkuka} again
\end{itemize}

\begin{tcolorbox}[breakable,colback=teal!4,colframe=teal!35!black,title={Before you move on}]
What does \ml{ഓണാക്കുക} mean? (\textquotedblleft{}To switch on\textquotedblright{}.) Say it once more.
\end{tcolorbox}
\section[\texorpdfstring{\ml{ഓഫാക്കുക} (ōphākkuka)}{ഓഫാക്കുക (ōphākkuka)}]{\ml{ഓഫാക്കുക} (ōphākkuka) — to switch off}

\label{lesson:ML-C188-ophakkuka}



\begin{quote}
Before the new one: say the Malayalam for to erase, then the Malayalam for to switch on.
\end{quote}

\subsection*{You'll want to know: \ml{ഓഫാക്കുക}}
\textbf{\ml{ഓഫാക്കുക}} — \emph{ōphākkuka} — \textquotedblleft{}to switch off\textquotedblright{}.

\begin{grammarlens}[title={What you've built}]
One more word for this chapter.
\end{grammarlens}

\subsection*{Your turn}
\begin{itemize}
\raggedright
  \item \emph{Say it:} \emph{ōphākkuka}
  \item \emph{Say it:} \emph{ōphākkuka}, once more
  \item \emph{Say it:} say \emph{ōphākkuka}
  \item \emph{Recall:} say the Malayalam for to erase, then the Malayalam for to switch on, then say \emph{ōphākkuka} again
\end{itemize}

\begin{tcolorbox}[breakable,colback=teal!4,colframe=teal!35!black,title={Before you move on}]
What does \ml{ഓഫാക്കുക} mean? (\textquotedblleft{}To switch off\textquotedblright{}.) Say it once more.
\end{tcolorbox}
\section[\texorpdfstring{\ml{പ്രസിദ്ധീകരിക്കുക} (prasiddhīkarikkuka)}{പ്രസിദ്ധീകരിക്കുക (prasiddhīkarikkuka)}]{\ml{പ്രസിദ്ധീകരിക്കുക} (prasiddhīkarikkuka) — to publish}

\label{lesson:ML-C188-prasiddhikarikkuka}



\begin{quote}
Before the new one: say the Malayalam for to switch on, then the Malayalam for to switch off.
\end{quote}

\subsection*{You'll want to know: \ml{പ്രസിദ്ധീകരിക്കുക}}
\textbf{\ml{പ്രസിദ്ധീകരിക്കുക}} — \emph{prasiddhīkarikkuka} — \textquotedblleft{}to publish\textquotedblright{}.

\begin{grammarlens}[title={What you've built}]
One more word for this chapter.
\end{grammarlens}

\subsection*{Your turn}
\begin{itemize}
\raggedright
  \item \emph{Say it:} \emph{prasiddhīkarikkuka}
  \item \emph{Say it:} \emph{prasiddhīkarikkuka}, once more
  \item \emph{Say it:} say \emph{prasiddhīkarikkuka}
  \item \emph{Recall:} say the Malayalam for to switch on, then the Malayalam for to switch off, then say \emph{prasiddhīkarikkuka} again
\end{itemize}

\begin{tcolorbox}[breakable,colback=teal!4,colframe=teal!35!black,title={Before you move on}]
What does \ml{പ്രസിദ്ധീകരിക്കുക} mean? (\textquotedblleft{}To publish\textquotedblright{}.) Say it once more.
\end{tcolorbox}
\section[\texorpdfstring{\ml{നന്നാക്കുക} (nannākkuka)}{നന്നാക്കുക (nannākkuka)}]{\ml{നന്നാക്കുക} (nannākkuka) — to repair, to fix}

\label{lesson:ML-C188-nannakkuka}



\begin{quote}
Before the new one: say the Malayalam for to switch off, then the Malayalam for to publish.
\end{quote}

\subsection*{You'll want to know: \ml{നന്നാക്കുക}}
\textbf{\ml{നന്നാക്കുക}} — \emph{nannākkuka} — \textquotedblleft{}to repair, to fix\textquotedblright{}.

\begin{grammarlens}[title={What you've built}]
One more word for this chapter.
\end{grammarlens}

\subsection*{Your turn}
\begin{itemize}
\raggedright
  \item \emph{Say it:} \emph{nannākkuka}
  \item \emph{Say it:} \emph{nannākkuka}, once more
  \item \emph{Say it:} say \emph{nannākkuka}
  \item \emph{Recall:} say the Malayalam for to switch off, then the Malayalam for to publish, then say \emph{nannākkuka} again
\end{itemize}

\begin{tcolorbox}[breakable,colback=teal!4,colframe=teal!35!black,title={Before you move on}]
What does \ml{നന്നാക്കുക} mean? (\textquotedblleft{}To repair, to fix\textquotedblright{}.) Say it once more.
\end{tcolorbox}
